\section{Reduction to certificates for ten vertices}
\label{sec:reduction}

\begin{theorem}[{\cite[Theorem~25(i) and Table~3]{firsching-2017}}]
  \label{thm:count}
  There are exactly $162004$ combinatorial types of simplicial $4$-polytopes
  with $10$ vertices.
\end{theorem}

Firsching's data \cite{firsching-2017} give a polytope with rational
vertices for each line of a file with $162004$ lines \artref{data}; we label
them \texttt{firsching-10-0} to \texttt{firsching-10-162003}. The proof of
the main result rests on the following two statements, which are finite
computations. The first is the analogue of
\cite[Lemma~6.1]{medgyes-maxvol-d4-n9} for ten vertices; the second is the
elimination of the other types.

\begin{hypothesis}
  \label{hyp:data}
  The polytopes \texttt{firsching-10-0} to \texttt{firsching-10-162003} are
  simplicial and pairwise not combinatorially equivalent, and for some
  $k_{0}$ the polytope \texttt{firsching-10-$k_{0}$} is combinatorially
  equivalent to the polytope of Lemma~\ref{lem:freesum} with $m = l = 5$.
\end{hypothesis}

\begin{hypothesis}
  \label{hyp:certificates}
  For every $k \neq k_{0}$ there are a coherent orientation $\mathcal{C}_{k}$
  of the facets of \texttt{firsching-10-$k$}, that is, a coherently oriented
  $\mathcal{C}_{k} = (10, \mathcal{F}_{k})$ whose facet vertex sets are
  exactly those of \texttt{firsching-10-$k$}, and a convexity certificate for
  $\mathcal{C}_{k}$ and some $\beta_{k} \in \mathbb{Q}$ with
  $\beta_{k} < \frac{25}{16}(10 + 2\sqrt{5})$.
\end{hypothesis}

By Theorem~\ref{thm:count}, Hypothesis~\ref{hyp:data} says that the
$162004$ polytopes represent every combinatorial type of simplicial
$4$-polytope with ten vertices, each exactly once. The polytopes need not be
inscribed: a convexity certificate is a polynomial identity and needs no
realisation, and for a type without an inscribed realisation the conclusion
of Proposition~\ref{prop:convexity-certificate} holds vacuously.
Hypothesis~\ref{hyp:certificates} allows $\lambda_{F,j} = 0$, so a
certificate in the sense of Definition~\ref{def:certificate} also serves.

\begin{remark}
  \label{rem:exact-bound}
  The condition on $\beta_{k}$ is decided exactly. With
  $t = \frac{16}{25}\beta_{k} - 10$, a rational number,
  $\beta_{k} < \frac{25}{16}(10 + 2\sqrt{5})$ if and only if $t < 2\sqrt{5}$,
  that is, if and only if $t < 0$ or $t^{2} < 20$.
\end{remark}

\begin{proposition}
  \label{prop:rivals}
  Assume Hypotheses~\ref{hyp:data} and~\ref{hyp:certificates}, and let
  $P \in \mathcal{P}_{4}(10)$ be simplicial. If $P$ is not combinatorially
  equivalent to \texttt{firsching-10-$k_{0}$}, then $\vol_{4}(P) < V^{*}$,
  with $V^{*}$ as in \eqref{eq:vstar}.
\end{proposition}

\begin{proof}
  By Hypothesis~\ref{hyp:data} and Theorem~\ref{thm:count}, $P$ is
  combinatorially equivalent to \texttt{firsching-10-$k$} for some
  $k \neq k_{0}$. Since $P$ is simplicial, a bijection between the vertex
  sets carries the facet vertex sets of this polytope, which are those of
  $\mathcal{C}_{k}$, onto those of $P$. So the vertices of $P$ can be listed
  as $p = (p_{1}, \ldots, p_{10})$ so that $p$ realises $\mathcal{C}_{k}$, and
  Proposition~\ref{prop:convexity-certificate} and
  Hypothesis~\ref{hyp:certificates} give
  \[
    \vol_{4}(P) \leq \frac{\beta_{k}}{24}
      < \frac{25}{384}\bigl(10 + 2\sqrt{5}\bigr) = V^{*}.
  \]
\end{proof}

\begin{theorem}
  \label{thm:main}
  Assume Hypotheses~\ref{hyp:data} and~\ref{hyp:certificates}. Then
  \[
    v_{4}(10) = \frac{25}{24}\sin^{2}\frac{2\pi}{5}
      = \frac{25}{384}\bigl(10 + 2\sqrt{5}\bigr) = 0.942196\ldots,
  \]
  and every $P \in \mathcal{P}_{4}(10)$ with $\vol_{4}(P) = v_{4}(10)$ is, up to
  an orthogonal transformation, the convex hull of two regular pentagons
  inscribed in $\mathbb{S}^{3}$ in orthogonal $2$-planes through the origin.
\end{theorem}

\begin{proof}
  By Lemma~\ref{lem:freesum} and Corollary~\ref{cor:pair} with
  $m = l = 5$, the pair of pentagons lies in $\mathcal{P}_{4}(10)$ and has
  volume $V^{*}$, so $v_{4}(10) \geq V^{*}$.

  Let $P \in \mathcal{P}_{4}(10)$ with $\vol_{4}(P) = v_{4}(10)$, which exists
  by Remark~\ref{rem:attained}. By Remark~\ref{rem:propertyz} and
  Lemma~\ref{lem:simplicial}, $P$ is simplicial. Since
  $\vol_{4}(P) \geq V^{*}$, Proposition~\ref{prop:rivals} shows that $P$ is
  combinatorially equivalent to \texttt{firsching-10-$k_{0}$}, and so, by
  Hypothesis~\ref{hyp:data}, to the polytope of Lemma~\ref{lem:freesum} with
  $m = l = 5$. Corollary~\ref{cor:pair} gives
  $v_{4}(10) = \vol_{4}(P) \leq V^{*}$, so $v_{4}(10) = V^{*}$, and its
  equality case identifies $P$.
\end{proof}

\begin{remark}
  \label{rem:computations}
  Hypothesis~\ref{hyp:data} is checked as for nine vertices
  \cite[Lemma~6.1]{medgyes-maxvol-d4-n9}: the facets of each polytope are
  computed from its rational vertices by exact supporting-hyperplane tests,
  each $3$-subset of a facet must lie in exactly two facets, and the
  $162004$ facet systems must have pairwise distinct canonical forms under
  relabelling. Hypothesis~\ref{hyp:certificates} is checked by the
  verification of Section~\ref{sec:verification} for each $k \neq k_{0}$,
  together with Remark~\ref{rem:exact-bound}. The certificates are found by
  semidefinite programming and rounded to rationals as in
  \cite[Section~3.4]{medgyes-maxvol-d4}; none of those steps enters the
  proof.
\end{remark}
